% !TeX root = closed-systems-from-comparison-completeness.tex
% ============================================================
\chapter{Comparison Triangles and the Multi--Point Boundary}
\label{ch:triangles}
% ============================================================
\section{Introduction}
Under the standing principle of closed-world admissibility
(\cref{sp:closed-world}), this chapter develops the first finite visible
transport obstruction in the sense of \cref{sp:transport}.
Proceeding from the lift classification, it identifies the first finite
triangle boundary that underlies \cref{ch:descent_obstruction,ch:bridge}.
\Cref{ch:locus,ch:lift} established that representative enrichment is
exhausted by representative choice and morphism-level transport,
possibly with both and with no third locus; equivalently, every
representative lift is given by representative choice together with
transport data, and no third layer of enrichment is present
(\cref{thm:lift-classification-main,cor:two-loci-lift}).
In particular, once a quotient-level protocol is lifted to
representatives in \(X\), the only possible obstruction to global
compatibility lies in the transport data.
The present chapter identifies the first nontrivial finite regime in
which that obstruction becomes intrinsically visible.
At the one-point level there is no issue: quotient semantics is complete,
and admissible reports are exactly the \(G\)-invariant maps
\[
	R:X\to S,
\]
equivalently those that factor through the orbit projection
\[
	\pi:X\to \Phys:=X/G
\]
(\cref{lem:forced-descent}).
At the multi-point level, however, a stronger compatibility problem
appears. Several representatives may each be individually compatible
with the quotient data, yet fail to admit a \emph{single common diagonal
	gauge} aligning them simultaneously.
This failure is the first genuinely relational boundary beyond ordinary
quotient semantics.
The chapter has three aims.
\begin{enumerate}
	\item We formulate diagonal \(n\)-point semantics and show that
	      \(n\)-point coherent observables separate exactly the diagonal
	      \(G\)-orbits in \(X^n\).
	\item We prove a stabilizer--coset criterion for simultaneous diagonal
	      alignment. This converts the problem into a concrete intersection
	      problem inside \(G\).
	\item We define the \emph{minimal detection size}
	      \[
		      m=\min\{|I|:\bigcap_{i\in I}C_i=\varnothing\},
	      \]
	      prove the associated minimality theorem, and isolate the regime \(m=3\),
	      called the \emph{triangle regime}. In that regime, pairwise compatibility
	      still holds but global compatibility fails, so triangles are the first
	      nontrivial finite witnesses of transport obstruction.
\end{enumerate}
This chapter is entirely algebraic. No topology, smoothness, metric,
probability, or dynamics is assumed.
% ============================================================
\section{\texorpdfstring{Diagonal \(n\)-point semantics}{Diagonal n-point semantics}}
% ============================================================
Let \(G\) act diagonally on \(X^n\) by
\[
	g\cdot(x_1,\dots,x_n):=(g\cdot x_1,\dots,g\cdot x_n).
\]
\begin{definition}[Diagonal orbit projection]
	For each \(n\ge 1\), let
	\[
		\pi_n:X^n\to X^n/G
	\]
	denote the orbit projection for the diagonal action.
\end{definition}
\begin{definition}[Diagonal equivalence]
	For \(x,y\in X^n\), write
	\[
		x\sim_\Delta y
	\]
	if there exists \(g\in G\) such that
	\[
		y=g\cdot x.
	\]
	Equivalently,
	\[
		x\sim_\Delta y
		\iff
		\pi_n(x)=\pi_n(y).
	\]
\end{definition}
\begin{definition}[\(n\)-point coherent observable]
	Let \(S\) be a set. A map
	\[
		R^{(n)}:X^n\to S
	\]
	is called \emph{\(n\)-point coherent} if it is invariant under the
	diagonal action of \(G\), i.e.
	\[
		R^{(n)}(g\cdot x)=R^{(n)}(x)
		\qquad
		\text{for all }g\in G,\ x\in X^n.
	\]
\end{definition}
\begin{theorem}[Diagonal separation]
	\label{thm:diag-separation}
	For \(x,y\in X^n\), the following are equivalent.
	\begin{enumerate}
		\item
		      \[
			      \pi_n(x)=\pi_n(y).
		      \]
		\item Every \(n\)-point coherent observable takes the same value on
		      \(x\) and \(y\).
	\end{enumerate}
\end{theorem}
\begin{proof}
	Assume first that \(\pi_n(x)=\pi_n(y)\). Then \(y=g\cdot x\) for some
	\(g\in G\). Hence for every \(n\)-point coherent observable
	\(R^{(n)}\),
	\[
		R^{(n)}(y)=R^{(n)}(g\cdot x)=R^{(n)}(x).
	\]
	Conversely, the quotient map
	\[
		\pi_n:X^n\to X^n/G
	\]
	is itself invariant under the diagonal action. Therefore, if every
	diagonally invariant map agrees on \(x\) and \(y\), then in particular
	\[
		\pi_n(x)=\pi_n(y).
	\]
\end{proof}
\begin{remark}[Multi-point boundary]
	For \(n=1\), \cref{thm:diag-separation} is just ordinary quotient
	semantics. For \(n\ge 2\), diagonal orbit structure on \(X^n\) may
	contain strictly more information than the pointwise orbit list in
	\(\Phys^n\), because diagonal equivalence requires one common element of
	\(G\) for all coordinates simultaneously.
\end{remark}
% ============================================================
\section{Alignment and stabilizer cosets}
% ============================================================
We now convert simultaneous diagonal alignment into an intersection
problem in the group \(G\).
\begin{definition}[Stabilizers and alignment cosets]
	\label{def:cosets}
	Fix \(n\ge 1\), points
	\[
		x_1,\dots,x_n\in X,
	\]
	and elements
	\[
		g_1,\dots,g_n\in G.
	\]
	Define
	\[
		y_k:=g_k^{-1}\cdot x_k
		\qquad
		(1\le k\le n).
	\]
	Let
	\[
		H_k:=\Stab(x_k)
		=
		\{h\in G:\ h\cdot x_k=x_k\}.
	\]
	Define the associated alignment coset
	\[
		C_k:=H_k g_k^{-1}\subseteq G.
	\]
\end{definition}
\begin{lemma}[Coset intersection criterion]
	\label{lem:coset}
	With notation as in \cref{def:cosets},
	\[
		(y_1,\dots,y_n)\sim_\Delta (x_1,\dots,x_n)
		\iff
		\bigcap_{k=1}^n C_k\neq\varnothing.
	\]
\end{lemma}
\begin{proof}
	By definition,
	\[
		(y_1,\dots,y_n)\sim_\Delta (x_1,\dots,x_n)
	\]
	if and only if there exists \(a\in G\) such that
	\[
		a\cdot x_k=y_k
		\qquad
		\text{for all }k.
	\]
	Using \(y_k=g_k^{-1}\cdot x_k\), this is equivalent to
	\[
		a\cdot x_k=g_k^{-1}\cdot x_k
		\qquad
		\text{for all }k.
	\]
	Multiplying on the left by \(g_k\) gives
	\[
		g_k a\cdot x_k=x_k
		\qquad
		\text{for all }k,
	\]
	hence
	\[
		g_k a\in H_k
		\qquad
		\text{for all }k.
	\]
	Equivalently,
	\[
		a\in H_k g_k^{-1}=C_k
		\qquad
		\text{for all }k.
	\]
	Thus such an \(a\) exists if and only if
	\[
		a\in\bigcap_{k=1}^n C_k.
	\]
\end{proof}
\begin{remark}[Free-action specialization]
	If each stabilizer \(H_k\) is trivial, then
	\[
		C_k=\{g_k^{-1}\},
	\]
	and \cref{lem:coset} reduces to the statement that simultaneous diagonal
	alignment holds if and only if
	\[
		g_1=\cdots=g_n.
	\]
	Thus in the free regime the multi-point obstruction is exactly the
	failure of all transport labels to coincide.
\end{remark}
% ============================================================
\section{Minimal detection size}
% ============================================================
\begin{definition}[Minimal detection size]
	With notation as in \cref{def:cosets}, define
	\[
		m
		:=
		\min
		\left\{
		|I|:
		I\subseteq\{1,\dots,n\},
		\ \bigcap_{i\in I}C_i=\varnothing
		\right\},
	\]
	with the convention \(m=\infty\) if no such subset exists.
\end{definition}
Thus \(m\) is the smallest number of points whose simultaneous diagonal
alignment fails.
\begin{remark}[Restriction notation]
	If \(I=\{i_1,\dots,i_r\}\subseteq\{1,\dots,n\}\), we write
	\[
		(x_i)_{i\in I}:=(x_{i_1},\dots,x_{i_r}),
		\qquad
		(y_i)_{i\in I}:=(y_{i_1},\dots,y_{i_r}),
	\]
	in the induced order on \(I\).
\end{remark}
\begin{theorem}[Minimal obstruction theorem]
	In the local alignment setup, assume \(m<\infty\). Then:
	\begin{enumerate}
		\item For every \(I\subseteq\{1,\dots,n\}\) with \(|I|<m\),
		      \[
			      (y_i)_{i\in I}\sim_\Delta (x_i)_{i\in I}.
		      \]
		\item There exists \(I\subseteq\{1,\dots,n\}\) with \(|I|=m\) such that
		      \[
			      (y_i)_{i\in I}\not\sim_\Delta (x_i)_{i\in I}.
		      \]
	\end{enumerate}
	Hence \(m\) is the unique minimal cardinality at which failure of a
	common diagonal gauge can be detected.
\end{theorem}
\begin{proof}
	Let \(I\subseteq\{1,\dots,n\}\) with \(|I|<m\). By minimality of \(m\),
	one has
	\[
		\bigcap_{i\in I}C_i\neq\varnothing.
	\]
	Applying \cref{lem:coset} to the restricted tuples gives
	\[
		(y_i)_{i\in I}\sim_\Delta (x_i)_{i\in I}.
	\]
	This proves the first statement.
	By definition of \(m\), there exists \(I\) with \(|I|=m\) and
	\[
		\bigcap_{i\in I}C_i=\varnothing.
	\]
	Applying \cref{lem:coset} again yields
	\[
		(y_i)_{i\in I}\not\sim_\Delta (x_i)_{i\in I}.
	\]
	This proves the second statement.
\end{proof}
\begin{corollary}[\(m=2\) and \(m=3\) regimes]
	With notation as in \cref{def:cosets}:
	\begin{enumerate}
		\item \(m=2\) if and only if there exist \(i\neq j\) such that
		      \[
			      C_i\cap C_j=\varnothing.
		      \]
		\item \(m=3\) if and only if every pairwise intersection is nonempty,
		      but there exist distinct \(i,j,k\) such that
		      \[
			      C_i\cap C_j\cap C_k=\varnothing.
		      \]
	\end{enumerate}
\end{corollary}
\begin{proof}
	This is immediate from the definition of \(m\).
\end{proof}
% ============================================================
\section{The triangle regime}
% ============================================================
\begin{definition}[Triangle regime]
	We say that the configuration is in the \emph{triangle regime} if
	\[
		m=3.
	\]
\end{definition}
\begin{theorem}[Strict minimality of triangles]
	\label{thm:triangle}
	In the local alignment setup, assume the triangle regime. Then there
	exist distinct indices
	\(i,j,k\) such that:
	\begin{enumerate}
		\item every two-point restriction among \(\{i,j,k\}\) is diagonally
		      alignable;
		\item the three-point restriction
		      \[
			      (y_i,y_j,y_k)
		      \]
		      is not diagonally alignable with
		      \[
			      (x_i,x_j,x_k).
		      \]
	\end{enumerate}
	Consequently:
	\begin{enumerate}
		\item no two-point coherent observable detects the obstruction;
		\item some three-point coherent observable does detect it.
	\end{enumerate}
\end{theorem}
\begin{proof}
	By definition of the triangle regime, there exists a triple
	\(\{i,j,k\}\) such that
	\[
		C_i\cap C_j\cap C_k=\varnothing.
	\]
	Again because \(m=3\), every subset of cardinality \(<3\) has nonempty
	intersection. In particular,
	\[
		C_i\cap C_j\neq\varnothing,
		\qquad
		C_i\cap C_k\neq\varnothing,
		\qquad
		C_j\cap C_k\neq\varnothing.
	\]
	By \cref{lem:coset}, each corresponding two-point restriction is
	diagonally equivalent. This proves the first statement.
	Applying \cref{lem:coset} to the triple gives
	\[
		(y_i,y_j,y_k)\not\sim_\Delta (x_i,x_j,x_k).
	\]
	This proves the second statement.
	The detectability statements follow from \cref{thm:diag-separation}:
	the two-point restrictions are diagonally equivalent and therefore
	cannot be separated by any two-point coherent observable, whereas the
	three-point restriction is not diagonally equivalent and is therefore
	separated by some three-point coherent observable. The separating
	three-point observable may be taken to be the orbit projection
	\[
		\pi_3:X^3\to X^3/G.
	\]
\end{proof}
\begin{corollary}[Triangle detectability]
	In the triangle regime, \(n=3\) is strictly minimal for detecting the
	obstruction to global diagonal alignment.
\end{corollary}
\begin{proof}
	This is just a restatement of \cref{thm:triangle}.
\end{proof}
\begin{remark}[Structural meaning of the triangle regime]
	The significance of the triangle regime is not that every obstruction
	must already occur at size \(3\). The significance is that \(m=3\) is
	the first regime in which global incompatibility can persist despite
	complete pairwise compatibility. Thus triangles are the first
	nontrivial finite witnesses of genuinely relational transport failure.
\end{remark}
% ============================================================
\section{Transport realization from lifts}
% ============================================================
We now connect the coset formalism of
\cref{def:cosets,lem:coset} directly to the lift and
transport structure of \cref{ch:lift}, without introducing any
auxiliary base object or path choice.
\begin{definition}[Finite transport sample]
	\label{def:transport-sample}
	Let
	\[
		\overline R:\mathcal I\to\Phys^{\disc}
	\]
	be a quotient-level protocol, and let
	\[
		T:\mathcal I\to X\sslash G
	\]
	be a representative lift.
	A \emph{finite transport sample} extracted from \(T\) consists of:
	\begin{enumerate}
		\item a finite family of objects
		      \[
			      i_1,\dots,i_n\in\Obj(\mathcal I),
		      \]
		\item representatives
		      \[
			      x_k:=T(i_k)\in X
			      \qquad
			      (1\le k\le n),
		      \]
		\item morphisms in the action groupoid
		      \[
			      (g_k,x_k):x_k\to y_k
			      \qquad
			      (1\le k\le n),
		      \]
		      arising from the transport data carried by \(T\).
	\end{enumerate}
	For such a sample, define
	\[
		H_k:=\Stab(x_k),
		\qquad
		C_k:=H_k g_k^{-1}.
	\]
\end{definition}
\begin{lemma}[Intrinsic character of the induced cosets]
	\label{lem:intrinsic-cosets}
	Let a finite transport sample be extracted from a representative lift
	\(T\). Then the associated cosets
	\[
		C_k=H_k g_k^{-1}
	\]
	depend only on the transport morphisms of the lift and are unchanged by
	replacing any chosen transport label \(g_k\) by another label
	representing the same morphism with the same source \(x_k\).
	Consequently all intersection properties of the family \(\{C_k\}\) are
	intrinsic to the lift.
\end{lemma}
\begin{proof}
	Fix \(k\). In the action groupoid, a morphism with source \(x_k\) and
	target \(y_k\) may be written as
	\[
		(g_k,x_k):x_k\to y_k.
	\]
	Suppose the same morphism is represented by another group element
	\(g_k'\in G\) with
	\[
		(g_k',x_k):x_k\to y_k.
	\]
	Then
	\[
		g_k\cdot x_k=y_k=g_k'\cdot x_k,
	\]
	hence
	\[
		g_k' g_k^{-1}\cdot x_k=x_k.
	\]
	Therefore
	\[
		g_k' g_k^{-1}\in H_k.
	\]
	Write
	\[
		g_k'=h_k g_k
		\qquad
		\text{for some } h_k\in H_k.
	\]
	Then
	\[
		H_k (g_k')^{-1}
		=
		H_k g_k^{-1} h_k^{-1}
		=
		H_k g_k^{-1},
	\]
	because \(h_k^{-1}\in H_k\).
	Thus the coset \(C_k\) is unchanged.
	Since this holds for each \(k\), all intersections of the family
	\(\{C_k\}\) are unchanged as well.
\end{proof}
\begin{proposition}[Transport induces alignment cosets]
	Every finite transport sample extracted from a representative lift
	determines canonically a family of alignment cosets
	\[
		C_k=H_k g_k^{-1},
		\qquad
		H_k=\Stab(x_k),
	\]
	and the associated alignment problem is exactly the intrinsic alignment
	problem of \cref{lem:coset}.
\end{proposition}
\begin{proof}
	This is immediate from \cref{def:transport-sample},
	\cref{lem:intrinsic-cosets}, and \cref{lem:coset}.
\end{proof}
\begin{corollary}[Triangle obstruction for transport]
	If the cosets induced by a finite transport sample lie in the triangle
	regime, then every pair of transported representatives is
	simultaneously alignable, but some triple is not.
	Thus the first nontrivial finite obstruction to global transport
	alignment is triangular.
\end{corollary}
\begin{proof}
	Immediate from \cref{thm:triangle}.
\end{proof}
\begin{remark}[Morphism-locus character of the obstruction]
	The obstruction detected here does not arise from quotient semantics
	itself. It arises only after representatives have been chosen and
	transport labels compared across several points. Thus the triangle
	regime lies strictly in the morphism locus of enrichment, i.e.\
	strictly beyond the object-level representative data isolated in
	\cref{cor:two-loci-lift,thm:flat-connection-interpretation}.
\end{remark}
% ============================================================
\section{Interpretation}
% ============================================================
The multi-point obstruction studied here is the first place where the
morphism-locus transport data of \cref{ch:lift} become visible as
an intrinsically relational phenomenon.
At one point, quotient semantics forgets all representative dependence.
At two points, there may or may not already be an obstruction. But in
the triangle regime, every pair remains compatible while the triple
fails. This is the first configuration in which pairwise transport
coherence ceases to imply global coherence.
The conclusion is exact: the triangle regime is the first regime in
which compatibility ceases to be controlled by lower-arity data. In that
sense, triangles are the first irreducible finite carriers of global
transport incompatibility.
This is the precise algebraic precursor of the later geometric picture.
The next chapter globalizes this triangular incompatibility into a
loop-based descent obstruction, and \cref{ch:bridge} then identifies
that loop defect and the triangle witness as two presentations of the
same intrinsic morphism-locus invariant. Under the subsequent
loop-level globalization of \cref{ch:descent_obstruction} and the
triangle--loop identification of \cref{ch:bridge}, this same finite
witness becomes the corresponding relational loop defect. Under smooth
realization, the infinitesimal form of that loop defect is curvature.

\section{Conclusion}
\label{sec:triangles-conclusion}
The finite obstruction boundary is now sharp: the triangle regime is the
first irreducible configuration in which transport data are pairwise
compatible yet globally inconsistent. No lower-arity test detects this
failure.

Accordingly, \cref{ch:descent_obstruction} globalizes this finite
witness into a loop-representation descent obstruction, and
\cref{ch:bridge} then identifies the finite triangle witness and loop
defect as two presentations of the same morphism-level obstruction
datum.
